В этой главе демонстрируются базовые средства оформления связного математического текста.
Сначала мы рассмотрим различие между алгебраическими выражениями, логическими выражениями и определениями, а затем покажем способы оформления выключных формул, ссылок, списков, рисунков и таблиц.
Раздел о математическом тексте имеет самостоятельную содержательную ценность и может использоваться как краткое введение в язык математических рассуждений.
Остальная часть главы служит демонстрацией типографских и структурных возможностей шаблона.